% Generated draft: requires mathematical and source review.
\documentclass[11pt]{article}
\usepackage[margin=25mm]{geometry}
\usepackage{fontspec}
\setmainfont{Latin Modern Roman}
\usepackage{amsmath,amssymb,mathtools,unicode-math}
\setmathfont{Latin Modern Math}
\usepackage{enumitem,needspace,xurl,hyperref}
\setlength{\parindent}{0pt}
\setlength{\parskip}{5pt}
\setlength{\emergencystretch}{3em}
\newcommand{\N}{\mathbb{N}}
\newcommand{\Z}{\mathbb{Z}}
\newcommand{\Q}{\mathbb{Q}}
\newcommand{\R}{\mathbb{R}}
\newcommand{\C}{\mathbb{C}}
\newcommand{\F}{\mathbb{F}}
\newcommand{\A}{\mathbb{A}}
\newcommand{\PP}{\mathbb{P}}
\newcommand{\E}{\mathbb{E}}
\newcommand{\eps}{\varepsilon}
\newcommand{\dd}{\,\mathrm{d}}
\begin{document}
\section*{UnsolvedMath ID: 1100501}
\textbf{Source problem code:} catalogue entry
\subsection{Definitions and mathematical statement}
Let $\mathbb H^n$ be real hyperbolic space and $G\le\operatorname{Isom}(\mathbb H^n)$. Under the standard axis convention set
\[
\operatorname{Ax}(G)=\{\operatorname{axis}(g):g\in G\text{ is loxodromic}\},
\]
where an axis is the unoriented geodesic line on which $g$ acts by nonzero translation, possibly with a rotational component in higher dimensions. Parabolic and elliptic elements contribute no such axes. Two subgroups in this fixed ambient group are commensurable if their intersection has finite index in each.

Anderson's literal source question is whether, for every pair of finitely generated discrete groups $G_1,G_2\le\operatorname{Isom}(\mathbb H^n)$,
\[
\operatorname{Ax}(G_1)=\operatorname{Ax}(G_2)
\quad\Longrightarrow\quad
[G_1:G_1\cap G_2]<\infty\ \text{ and }\ [G_2:G_1\cap G_2]<\infty.
\]
The source gives no range of $n$; the usual setting is $n\ge2$. The supplied later research summary warns that the literal assertion is false for elementary parabolic groups, which have empty axis sets, and identifies the intended open question as the same implication with both groups non-elementary. In the usual convention non-elementary means that the limit set in $\partial\mathbb H^n$ has more than two points. The non-elementary restriction is an explicit source-supplied correction, not a hypothesis present in the original displayed question.

For context, the summary proposes studying the axis-set stabilizer and commensurator,
\[
\Sigma(G)=\{h\in\operatorname{Isom}(\mathbb H^n):h\operatorname{Ax}(G)=\operatorname{Ax}(G)\},\qquad
\operatorname{Comm}(G)=\{h:hGh^{-1}\text{ is commensurable with }G\},
\]
and the equality $\Sigma(G)=\operatorname{Comm}(G)$. No precise reduction theorem or hypotheses establishing equivalence with the question are supplied, so this is recorded as a suggested framework rather than an asserted equivalent formulation.

The supplied summary reports affirmative results when both groups are arithmetic Fuchsian, or both are arithmetic Kleinian, and partial results for other Kleinian classes. It describes the general non-arithmetic non-elementary problem as open, including the case of Fuchsian groups isoaxial with $\operatorname{PSL}(2,\Z)$, where isoaxial means equality of the axis sets. These are supplied status claims, not independently verified results.
\subsection{Short English statement}
If two finitely generated discrete hyperbolic isometry groups have exactly the same loxodromic axes, must they share a subgroup of finite index in each? The later source summary says a non-elementary hypothesis is needed and reports arithmetic Fuchsian and Kleinian cases, while leaving the general non-arithmetic question open. The original statement and this explicitly qualified intended version are kept distinct.
\subsection{Sources}
\begin{itemize}
\item Local source registry: \url{data/unsolvedmath/problems.json}, permanent UnsolvedMath ID $1100501$, source code catalogue entry, including the supplied OPEN-TRIAGE research summary. Dataset attribution: Ulam AI and UnsolvedMath Contributors; source attribution and rights are retained.
\item Mladen Bestvina, \emph{Questions in Geometric Group Theory}, source-list date 2004, Question 5.1, attributed to Jim Anderson; supplied metadata identifies PDF page 12: \url{https://www.math.utah.edu/~bestvina/eprints/questions-updated.pdf}. The metadata reports manual transcription and access on July 29, 2026, with current-status and rights review flags.
\item Long--Reid (1998), cited in the supplied summary for arithmetic Fuchsian groups; Reid's 2014 survey, Theorem 6.1, Question 7.1, and the remark after Theorem 6.1, cited for arithmetic Kleinian results and open non-arithmetic cases; Xie--Jiang (2008), cited for additional Kleinian partial results. The record supplies no complete titles or separate links.
\item A Mess (1990) preprint is mentioned by the supplied summary as motivating context. The summary states it was never published and its result is not independently verifiable; no precise result or complete citation is supplied.
\end{itemize}
\end{document}
